\section{Introduction}
\label{sec:introduction}

The philosophy of Heraclitus is often distilled into the
aphorism \textit{panta rhei}---everything flows.
In the physical world, appearances and spatial configurations
continually change, while objects maintain continuity over time.
An embodied world model should capture both aspects.
Recent generative video models~\citep{wan2025wan,seedance2026seedance2}
predict increasingly coherent visual futures, but express
these predictions in pixels, leaving the underlying 3D
dynamics implicit.
World action models~\citep{zhou2026zerowam,yuan2026fastwam,nvidia2026cosmos3,ye2026dreamzero}
connect future prediction to robotic behavior by incorporating
actions into world modeling.
However, action supervision is costly to collect at scale,
and action representations are often tied to particular
robot embodiments, complicating cross-platform transfer.
These limitations motivate extending world modeling into
the space of explicit 3D motion flow: a geometric representation
of physical change that can connect visual prediction to
planning and control across embodiments.

A complementary line of research approaches this space
through motion modeling, progressing from tracking motion
in observed videos~\citep{lu2026track4world,luo20264rc} to predicting future trajectories
from current observations~\citep{zhang2026molmomotion,lee2026mu_0}.
For generative motion prediction, producing dense 3D
scene flow is computationally expensive.
Existing scene-level methods instead predict motion
on coarse spatial grids~\citep{lee2026tracegen,li2026egowam}.
However, sparse scene flow still allocates motion tokens
to task-irrelevant regions, while its reduced spatial
resolution limits the fidelity of fine motion boundaries.
Query-conditioned prediction offers a more targeted
alternative, allocating motion tokens to specified 3D
locations in task-critical regions while retaining global
scene context as conditioning information.
Yet learning generalizable dynamics at scale faces two
challenges: high-quality 3D trajectory annotations are costly
to acquire, while flow-only supervision provides a limited
learning signal for physical dynamics, potentially limiting
generalization across scenes and interactions.
Jointly modeling visual evolution and explicit motion
offers a natural way to enrich this supervision,
bringing motion prediction into a generative world model.

To this end, we introduce \method{}, a generative world model
that jointly predicts future video and query-conditioned
3D object trajectories through a unified
Mixture-of-Transformers (MoT) backbone.
Given an observed image, a language instruction, and sparse
3D query points, \method{} forecasts motion at the queried
locations while using the full scene as visual context.
We represent trajectories as continuous tokens encoding
normalized frame-to-frame 3D displacements, enriched with
initial point geometry and spatiotemporally aligned with
video tokens.
A unified rectified-flow objective jointly denoises both
modalities, coupling trajectory learning with dense visual
supervision of the same future.

To support training at scale, we assemble a heterogeneous
corpus spanning simulated and real-world domains, comprising
approximately 1.44 million object-level samples---equivalent
to 1,224 hours of standardized 49-frame sequences at 16 FPS.
We combine existing motion annotations with automated
trajectory extraction, unify their representations, and
apply source-specific quality control and fixed domain-level
normalization for mixed-data training.
We further introduce WorldFlowBench, comprising three
evaluation sets, including SpatialCog for spatial reasoning
in future motion prediction.
On WorldFlowBench, \method{} achieves the lowest ADE and FDE
across all three evaluation sets among non-static methods.
Compared with its flow-only variant, joint video-flow
generation reduces ADE by 48.5\% on Sim7 and 55.9\% on
SpatialCog, with a 68.5\% reduction in FDE on SpatialCog.
These results support the value of complementary visual
supervision for accurate 3D motion forecasting.

Our contributions are:
\begin{itemize}
    \item We introduce \method{}, a generative world model
    that jointly predicts video and query-conditioned 3D
    object motion through a unified MoT backbone, combining
    explicit trajectory supervision with dense visual
    supervision.

    \item We construct a heterogeneous corpus of approximately
    1.44 million object-level samples from simulated and
    real-world data, using unified motion representations
    and source-specific quality control to support
    large-scale training.

    \item We introduce WorldFlowBench, comprising three
    evaluation sets including SpatialCog for spatial
    reasoning, and demonstrate improvements over existing
    motion prediction methods and
    video-generation-then-tracking pipelines.
\end{itemize}